% !TEX program = xelatex
% Черновик для задания по научному письму. Эксперименты ещё не проведены.
% Класс misisarticle.cls берётся из форка шаблона Садекова без изменений.
\documentclass{misisarticle}

\grnti{28.23.15}
\asjc{1707}
\oecd{1.2.01}

\title{Сегментация трехмерного облака точек нескольких ToF-датчиков в задаче измерения габаритов объектов}

\misisauthor{В.~А.~Охота}{кафедра инженерной кибернетики НИТУ «МИСиС»}{Москва, Россия}{m2306180@edu.misis.ru}
\misisauthor{С.~С.~Мишуров}{кафедра инженерной кибернетики НИТУ «МИСиС»}{Москва, Россия}{}

\abstract{%
Рассматривается задача выделения точек измеряемого объекта в трехмерной сцене, сформированной по данным нескольких ToF-датчиков. Цель исследования состоит в сравнении методов сегментации по качеству выделения объекта и погрешности определения его длины, ширины и высоты. Облака отдельных ракурсов предполагается преобразовывать в общую систему координат с использованием внешних параметров калибровки. Рассматриваются непосредственная обработка трехмерной геометрии и перенос двумерных масок на облако точек при наличии согласованных изображений. Сформулированы предварительная постановка задачи, требования к собственному набору данных и протокол оценки. Для сегментации предусматриваются метрики mIoU, precision, recall и F1-score, для габаритов --- абсолютная и относительная ошибки относительно контрольных измерений. Проверку методов планируется провести на собственном наборе многовидовых ToF-сцен упакованных грузов. Состав, объем и характеристики датасета, формат разметки и окончательные параметры эксперимента будут уточнены на этапе подготовки данных; собственные экспериментальные результаты в настоящем тексте не представлены.
}

\keywords{сегментация изображений, облака точек, SAM~2, проекция, калибровка камеры, классификация точек, ToF-датчики}

\begin{document}
\maketitle

\section{Введение}

Автоматическое определение габаритов применяется при учете, размещении и перемещении объектов. Для получения метрических размеров требуется отделить измеряемую поверхность от фона и связать наблюдения с геометрией сцены. ToF-датчики предоставляют информацию о расстоянии, по которой строятся облака точек. Однако шум, перекрытия и неполная наблюдаемость поверхности затрудняют измерение. Значимость этих факторов для определения размеров показана в работе Rodriguez и соавторов~\cite{rodriguez2023}.

Использование нескольких ракурсов позволяет наблюдать разные участки поверхности, но вводит дополнительный этап преобразования координат. В объединенное облако попадают как точки объекта, так и точки окружающих поверхностей. Поэтому качество измерения зависит от двух связанных процессов: согласования геометрии датчиков и определения принадлежности точек объекту. Изолированная оценка сегментации не описывает полностью точность конечного измерения.

Например, небольшое число ошибочно включенных удаленных точек может увеличить ограничивающий объем объекта, хотя пересечение сегмента с эталонной разметкой изменится незначительно. Потеря точек у крайней границы, напротив, способна уменьшить оценку размера. Следовательно, метод, дающий более высокий IoU, необходимо дополнительно проверять по погрешности габаритов. Это положение рассматривается как основание будущего сравнения, а не как установленный экспериментальный результат.

Цель исследования --- сопоставить методы сегментации объединенного ToF-облака по точности выделения объекта и определения его габаритов. Для достижения цели необходимо описать геометрические преобразования между датчиками, определить рассматриваемые способы сегментации, задать единый способ вычисления размеров и подготовить протокол проверки на собственных данных. Предполагаемый вклад состоит в выявлении влияния ошибок сегментации и объединения ракурсов на погрешность измерения габаритов и формулировании рекомендаций по выбору метода и его параметров.

На текущем этапе сформулированы постановка задачи и методика сравнения. Датасет находится на этапе планирования; сведения о его объеме, конфигурации датчиков и разметке пока не установлены. Поэтому конкретные алгоритмы, параметры и выводы об их эффективности будут уточняться по доступным данным.

\section{Анализ работ}

\subsection{Измерение размеров по данным глубины}

Rodriguez и соавторы рассматривают измерение прямоугольных и цилиндрических объектов по шумным, частично наблюдаемым ToF-данным~\cite{rodriguez2023}. В основе подхода лежит подгонка суперквадрик с ограничениями, учитывающими неполноту облака. Авторы анализируют влияние ориентации объекта и свойств опорной поверхности. Работа непосредственно связана с задачей измерения габаритов, однако ее результаты не являются проверкой сегментации сцены, объединенной из нескольких ToF-датчиков.

В студенческой работе Прохорова объединяются сегментация изображений, монокулярная оценка глубины и вычисление метрических размеров с использованием параметров камеры~\cite{prokhorov2026}. Рассматриваются модели Depth Anything V2 и ZoeDepth. Эта работа показывает прикладную связь между выделением объекта и измерением его размеров. Существенное отличие от рассматриваемой постановки состоит в источнике геометрии: глубина предсказывается по изображению, тогда как в настоящем исследовании предполагаются измеренные ToF-данные.

\subsection{Получение трехмерных сегментов по двумерным маскам}

SAI3D объединяет геометрическое разбиение трехмерной сцены и маски Segment Anything Model, полученные для разных изображений~\cite{yin2024}. Геометрические примитивы последовательно сливаются в сегменты экземпляров с учетом многовидовой согласованности. Метод важен для рассматриваемой темы как пример связи трехмерной геометрии и двумерной сегментации. При этом результаты на интерьерных сценах требуют отдельной проверки для условий измерительной установки.

В Any3DIS используются отслеживание двумерных масок с SAM~2, геометрические сведения сцены и оптимизация трехмерных предложений~\cite{nguyen2025}. В отличие от независимого выделения масок для каждого ракурса, отслеживание поддерживает связь наблюдений одного объекта. Перенос этого подхода на ToF-установку зависит от наличия согласованных изображений и пригодности наблюдений для отслеживания; измерение трехмерных точек само по себе не гарантирует наличие RGB-канала.

\subsection{Оценка сегментации и геометрического согласования}

Андронов сравнивает FCN и DeepLabV3 для семантической сегментации уличных изображений~\cite{andronov2025}. Работа относится к двумерным данным и полезна для обсуждения качества масок и сравнения сегментаторов. Ее выводы не переносятся непосредственно на трехмерные облака, однако оценка масок становится частью общего протокола, если они используются для назначения меток ToF-точкам.

Али и Садеков рассматривают оценку точности позиционирования по сопоставлению сегментированных изображений с цифровой картой~\cite{ali2025}. В статье сравниваются геометрически связанные области и линии. Хотя ее приложение отличается от измерения габаритов, она демонстрирует необходимость оценивать сегментацию совместно с геометрией и конечной прикладной ошибкой.

Рассмотренные источники охватывают измерение размеров, получение трехмерных сегментов и оценку геометрического согласования. В настоящем исследовании эти направления связываются в единой проверке: для каждого результата сегментации вычисляются и точечные метрики, и ошибки габаритов. Преимущество конкретного метода заранее не предполагается.

\section{Постановка задачи}

\subsection{Входные данные и допущения}

Пусть наблюдение сцены получено с $M$ датчиков. Для датчика $i$ задано облако $\mathcal{P}_i$ в его локальной системе координат и внешние параметры $R_i$, $t_i$, переводящие точки в общую систему сцены. В предварительной постановке объект считается неподвижным в течение получения набора ракурсов. Координаты и векторы переноса должны быть выражены в согласованных единицах измерения.

Для точки $\mathbf{x}_i$ переход в общую систему имеет вид
\begin{equation}
    \mathbf{x}=R_i\mathbf{x}_i+t_i.
    \label{eq:transform}
\end{equation}
Требуется проверить направление калибровочного преобразования: параметры обратного перехода нельзя использовать в формуле~\eqref{eq:transform} без обращения. Объединенное облако состоит из преобразованных наблюдений всех датчиков.

Если исходные данные представлены картами глубины, для восстановления координат потребуются внутренние параметры камер и описание формата глубины. Если используются двумерные маски, дополнительно необходимы изображения, геометрически согласованные с измеренной глубиной. Конкретные параметры установки будут зафиксированы после определения датасета.

\subsection{Требуемый результат}

Для выбранного измеряемого объекта необходимо определить множество принадлежащих ему точек $\widehat{\mathcal{S}}$ и вектор габаритов $\widehat{\mathbf{d}}=(\widehat{L},\widehat{W},\widehat{H})$. При оценке выделения выбранного объекта используется бинарная принадлежность точки: объект или фон. Одновременное измерение нескольких объектов потребует также согласования идентификаторов экземпляров между ракурсами.

Задача сегментации формулируется как построение функции $\widehat{y}(\mathbf{x})\in\{0,1\}$ для точек объединенного облака. Оценка габаритов выполняется по точкам с меткой объекта. Эталонные метки $y(\mathbf{x})$ и контрольные размеры $L^*$, $W^*$, $H^*$ предназначены для проверки результата и не должны подменять входные данные алгоритма.

\subsection{Планируемый набор данных}

Планируется формирование собственного датасета, включающего облака точек объектов, полученные с нескольких ToF-датчиков. В качестве объектов исследования предполагается использовать упакованные грузы. Состав, объем и характеристики данных, а также формат разметки будут уточнены на этапе подготовки экспериментов.

Для воспроизводимой проверки потребуется сохранять принадлежность наблюдений датчикам, параметры калибровки и единицы измерения. Для расчета качества сегментации необходима эталонная разметка точек, а для оценки габаритов --- независимые контрольные измерения размеров. При использовании масок дополнительно сохраняются соответствующие изображения и сведения об их согласовании с глубиной.

Объем выборки, число ракурсов и способы получения разметки пока не фиксируются. Если метод требует обучения или настройки по размеченным данным, контрольные сцены необходимо отделить от данных настройки. Разделение следует выполнять по объектам или сценам, исключая попадание близких наблюдений одного объекта одновременно в настройку и итоговую проверку.

\section{Описание методов}

\subsection{Подготовка и объединение наблюдений}

Предполагаемая обработка начинается с проверки корректности измерений и удаления недопустимых значений глубины. Далее облака преобразуются согласно формуле~\eqref{eq:transform}. Параметры предварительной фильтрации должны быть одинаковыми для сравниваемых сегментаторов либо их различие должно учитываться отдельно.

Плотность точек может различаться между датчиками и областями объекта. Поэтому при необходимости используется единый способ пространственного прореживания. Его масштаб выбирается с учетом минимальных измеряемых элементов: чрезмерное уменьшение плотности способно убрать точки крайних границ. При разметке и расчете метрик необходимо учитывать соответствие между исходными и обработанными точками.

\subsection{Геометрический способ выделения объекта}

В качестве базового подхода предполагается непосредственная обработка трехмерной геометрии. Для установки с выраженной опорной плоскостью могут использоваться выделение этой плоскости, удаление ее точек и группировка оставшихся точек по пространственной близости. Выбор измеряемой группы выполняется по заранее заданному правилу, например по положению объекта в рабочей зоне.

Применимость такого способа зависит от сцены. Контакт объекта с соседними предметами, малая дистанция до фона и недостаточная наблюдаемость поверхности могут затруднять разделение. Поэтому конкретные алгоритмы выделения плоскости и группировки, а также их пороги будут определены после получения данных. Геометрический способ рассматривается как воспроизводимый базовый вариант для сравнения, а не как заранее выбранное окончательное решение.

\subsection{Перенос двумерных масок в облако точек}

Второе направление предполагает получение масок на согласованных изображениях и назначение меток трехмерным точкам. Концептуальная основа такого переноса представлена в SAI3D и Any3DIS~\cite{yin2024,nguyen2025}. Использование SAM~2 возможно при наличии подходящих изображений и условий для задания подсказок или отслеживания.

Для чтения маски точка сцены сначала переводится в систему соответствующей камеры:
\begin{equation}
    \mathbf{x}_i=R_i^{\mathsf{T}}(\mathbf{x}-t_i).
\end{equation}
Для координат $(X_i,Y_i,Z_i)$ недисторсированная проекция определяется выражениями
\begin{equation}
\begin{aligned}
    u&=f_x X_i/Z_i+c_x,\\
    v&=f_y Y_i/Z_i+c_y,
\end{aligned}
\label{eq:projection}
\end{equation}
где $f_x$, $f_y$, $c_x$, $c_y$ --- внутренние параметры камеры. Эта запись предполагает корректный учет дисторсии и соглашения о глубине в используемых данных.

Метка назначается только точкам, расположенным перед камерой, попадающим в изображение и видимым в данном ракурсе. Проверка видимости должна сопоставлять глубину проекции с наблюдаемой поверхностью. Иначе маска переднего объекта может быть ошибочно перенесена на скрытые точки фона.

При объединении меток нескольких ракурсов требуется единое правило обработки противоречий. В предварительном варианте можно использовать голосование только по видимым наблюдениям. Отсутствие точки в поле зрения отдельного датчика не считается наблюдением фона. Способ объединения масок и его параметры должны проверяться отдельно от качества двумерного сегментатора.

\subsection{Вычисление габаритов}

Для сравнения сегментации предполагается использовать единый способ вычисления размеров. Пусть $\mathbf{q}(\mathbf{x})$ --- координаты точки в выбранной системе осей объекта. Ориентация осей задается геометрией установки либо оценивается одинаковым способом для всех методов. Применение осей сцены к произвольно повернутому объекту без такого согласования изменяет измеряемые величины.

Для выделенного множества точек габарит по оси $j$ определяется как
\begin{equation}
    \widehat{d}_j=
    \max_{\mathbf{x}\in\widehat{\mathcal{S}}}q_j(\mathbf{x})-
    \min_{\mathbf{x}\in\widehat{\mathcal{S}}}q_j(\mathbf{x}).
    \label{eq:dimensions}
\end{equation}
Правило соответствия осей длине, ширине и высоте фиксируется для всей проверки. Если вводится дополнительная фильтрация выбросов или используется подгонка геометрической модели, эта процедура должна применяться согласованно и описываться как отдельный этап измерения.

Формула~\eqref{eq:dimensions} дает размеры ограничивающего объема наблюдаемого сегмента. При неполном покрытии поверхности эти значения могут отличаться от полных физических размеров объекта. Следовательно, при проверке необходимо учитывать видимость крайних участков и полноту наблюдений; добавление ракурсов не гарантирует полного покрытия автоматически.

\subsection{Критерии сравнения}

Качество сегментации предполагается оценивать по IoU, precision, recall и F1-score. Для выбранного объекта пересечение с эталонным множеством $\mathcal{S}^*$ имеет вид
\begin{equation}
    IoU=\frac{|\widehat{\mathcal{S}}\cap\mathcal{S}^*|}
                  {|\widehat{\mathcal{S}}\cup\mathcal{S}^*|}.
\end{equation}
При бинарной постановке mIoU усредняет IoU классов объекта и фона; дополнительно следует отдельно приводить IoU объекта. Precision, recall и F1-score характеризуют ошибочное включение фоновых точек и потерю точек объекта. Для интерпретации результатов важно фиксировать способ усреднения метрик по сценам.

Для каждой компоненты габаритов вычисляются абсолютная и относительная ошибки:
\begin{equation}
\begin{aligned}
    e_j&=|\widehat{d}_j-d_j^*|,\\
    e_j^{\mathrm{rel}}&=\frac{|\widehat{d}_j-d_j^*|}{d_j^*}.
\end{aligned}
\end{equation}
Контрольные размеры $d_j^*$ должны быть получены независимо от проверяемого алгоритма. При представлении ошибок необходимо указывать единицы измерения и точность контрольных измерений.

Сначала методы предполагается сравнивать при одинаковых наблюдениях и фиксированной калибровке. Затем может отдельно оцениваться влияние числа ракурсов, фильтрации и ошибок преобразования координат. Такая организация позволяет различать влияние сегментации и геометрического согласования. Диапазоны параметров и число проверяемых сцен будут определены после подготовки данных.

Таким образом, протокол предусматривает совместный анализ точечных метрик и ошибок длины, ширины и высоты для одних и тех же объектов. Итоговые рекомендации по выбору метода должны основываться на будущей проверке, включая случаи хорошего IoU при заметной ошибке габаритов и случаи неполного наблюдения границ.

\begin{thebibliography}{9}

\bibitem{rodriguez2023}
Rodriguez B., Rangarajan P., Zhang X., Rajan D.
Dimensioning Cuboid and Cylindrical Objects Using Only Noisy and Partially Observed Time-of-Flight Data // Sensors. --- 2023. --- Vol.~23, No.~21. --- Art.~8673.
--- URL: \url{https://doi.org/10.3390/s23218673}.

\bibitem{prokhorov2026}
Прохоров Д.~Д. Определение геометрических размеров объектов на изображениях с использованием монокулярной оценки глубины // Искусственный интеллект в промышленных, коммерческих, медицинских и финансовых приложениях: сборник статей научно-технического семинара. --- М.: НИТУ «МИСИС», 2026. --- Вып.~5. --- С.~98--103.
--- URL: \url{https://misis.ru/files/-/41f69c37b68def41ea566922c56b898e/kik_2025_2.pdf}.

\bibitem{yin2024}
Yin Y., Liu Y., Xiao Y., Cohen-Or D., Huang J., Chen B.
SAI3D: Segment Any Instance in 3D Scenes // Proceedings of the IEEE/CVF Conference on Computer Vision and Pattern Recognition (CVPR). --- 2024. --- P.~3292--3302.
--- URL: \url{https://doi.org/10.1109/CVPR52733.2024.00317}.

\bibitem{nguyen2025}
Nguyen P., Luu M., Tran A., Pham C., Nguyen K.
Any3DIS: Class-Agnostic 3D Instance Segmentation by 2D Mask Tracking // Proceedings of the IEEE/CVF Conference on Computer Vision and Pattern Recognition (CVPR). --- 2025. --- P.~3636--3645.
--- URL: \url{https://openaccess.thecvf.com/content/CVPR2025/html/Nguyen_Any3DIS_Class-Agnostic_3D_Instance_Segmentation_by_2D_Mask_Tracking_CVPR_2025_paper.html}.

\bibitem{andronov2025}
Андронов К.~В. Сравнение и анализ моделей FCN и DeepLabV3 в задаче семантической сегментации объектов городской улицы // Искусственный интеллект в промышленных, коммерческих, медицинских и финансовых приложениях: сборник статей научно-технического семинара. --- М.: НИТУ «МИСИС», 2025. --- Вып.~3. --- С.~9--13.
--- URL: \url{https://misis.ru/files/-/7030ef7ab572d5227e50a55006446a76/kik_2024_2.pdf}.

\bibitem{ali2025}
Али Б., Садеков Р.~Н. Оценка точности позиционирования трамвая в кривых на основе данных карты и сегментированных изображений // Научно-технический вестник информационных технологий, механики и оптики. --- 2025. --- Т.~25, №~4. --- С.~771--779.
--- URL: \url{https://doi.org/10.17586/2226-1494-2025-25-4-771-779}.

\end{thebibliography}
\end{document}
